% TOP_PROBLEM: {"problemId": "problem.cramer-conjecture-prime-gaps", "canonicalTitle": "Cramer Conjecture on the Limsup of Normalized Prime Gaps", "releaseRank": 160, "primaryDomain": "number_theory_arithmetic_geometry", "primaryDomainLabel": "Number theory & arithmetic geometry", "displayStatus": "open", "releaseStatus": "open"}
% !TeX program = lualatex
% Definition and sources only; this file is not an LLM solution attempt.
\documentclass[11pt]{article}
\usepackage[margin=25mm]{geometry}
\usepackage{fontspec}
\setmainfont{Latin Modern Roman}
\usepackage{amsmath,amssymb,mathtools}
\usepackage{unicode-math}
\setmathfont{Latin Modern Math}
\usepackage{xurl,hyperref}
\hypersetup{colorlinks=true,urlcolor=blue}
\setcounter{secnumdepth}{0}
\setlength{\parindent}{0pt}
\setlength{\parskip}{5pt}
\setlength{\emergencystretch}{3em}
\begin{document}
\section*{\texorpdfstring{Cramer Conjecture on the Limsup of Normalized Prime Gaps}{Cramer Conjecture on the Limsup of Normalized Prime Gaps}}\label{160}
\subsection{Definitions and mathematical statement}
Let $p_n$ be the $n$th prime and $g_n=p_{n+1}-p_n$. With natural logarithms, the selected assertion is the exact constant statement
\[
 \limsup_{n\to\infty}\frac{g_n}{(\log p_n)^2}=1.
\]
Equivalently, for every ${\varepsilon}>0$, all sufficiently large $n$ satisfy
$g_n\le(1+{\varepsilon})(\log p_n)^2$, while infinitely many satisfy
$g_n\ge(1-{\varepsilon})(\log p_n)^2$. The equality of the limsup is materially stronger than an unspecified upper bound $g_n=O((\log p_n)^2)$. This section records the catalogue's precise constant-one version; the title alone should not be used to interchange different prime-gap conjectures or heuristics.
\par\smallskip\textit{Formulation and concept references:} \href{https://mathworld.wolfram.com/CramerConjecture.html}{[S1]}.

\subsection{Short English statement}
Is the largest recurring scale of consecutive prime gaps exactly the square of the logarithm of the smaller prime, with limiting upper constant one?
\subsection{Sources}
\begin{itemize}
\item[S1]Weisstein, E. W., Cramer Conjecture, From MathWorld - A Wolfram Resource, updated 16 August 2026.
\url{https://mathworld.wolfram.com/CramerConjecture.html}.
\textit{Formulation source.}
\end{itemize}
\end{document}
